\chapter{Fazit}\label{ch:fazit}

In dieser Arbeit konnte erfolgreich aufgezeigt werden, dass die Formalisierung des \Abk{DSA}-Regelwerks zur Charakter-Validierung möglich ist.
Es hat sich dabei gezeigt, dass die Formalisierung mit dem wissenschaftlich erprobten \Abk{ASP} unkompliziert war und leicht verständlich ist, sowohl bei der Formalisierung selbst als auch bei der Integration des genutzten Frameworks \Begriff{clingo}.
Aufgrund des Umfangs konnten nicht sämtliche relevanten Regeln und Charakter-Merkmale beachtet werden, sodass erst in Zukunft bewertet werden kann, ob nach vollständiger Formalisierung dies so bleibt.
Üblicherweise wird mit \Abk{ASP} eine Aussage über die Erfüllbarkeit eines Problems und mögliche Lösungen getroffen.
Dieses Vorgehen wurde hier \textquote{gebrochen}, da (bisher) jedes erstelle \Abk{LP} erfüllbar sein musste und nicht mögliche Charaktere (Merkmals-Kombinationen) gefunden werden sollten, sondern vordefinierte abzuleitende Fakten als Validierungsergebnis generiert werden mussten.

Das Problem hoher Migrationsaufwände der bisherigen \Abk{DSA}-Hilfsmittel auf die aktuelle 5te-Ausgabe, aufgrund einer (angenommenen) starken Verflechtung der fachlichen Logik (dem Regelwerk) mit der technischen Implementierung oder zu komplexen und schlecht wartbaren Software, wurde in dieser Arbeit beachtet.
Mit den im Konzept aufgestellten Konventionen, als auch den getroffenen Architektur- und Design-Entscheidungen sowie den qualitätssichernden Maßnahmen in der Realisierung wurde eine hohe Wartbarkeit der Software erreicht.
Auch die Wahl von Python als Programmiersprache war dafür von Vorteil.
Ebenfalls konnte in der Realisierung, durch einen modularen Aufbau der Regelbücher-\AbkP{LP}, die Umsetzung weiterer Regelbücher und Regelbuch-Erweiterungen berücksichtigt werden.
Für die Implementierung anderer Regelwerksausgaben bedarf es jedoch geringfügiger Anpassungen, sofern die Charakter-Modellierung abweicht, wie \zB durch neue Merkmale.
Die entwickelte \Code{Engine} mit dem generischen Ablauf der Charakter-Validierung, klaren Ein- (Charaktermodell) und Ausgabe-Formaten (Ergebnis-Fakten) sowie den ermöglichten Freiheiten bei der strukturellen Umsetzung der \AbkP{LP} stellt die zentralste Komponente dar.
Die entstandene Software selbst ist insgesamt eine solide Grundbasis für ein langfristiges quelloffenes Projekt.
Dies war neben der Formalisierung ein weiteres Ziel der Arbeit.
